\documentclass[11pt,a4paper]{article}
\usepackage[utf8]{inputenc}\usepackage[T1]{fontenc}
\usepackage[margin=2.3cm]{geometry}
\usepackage{amsmath,amssymb,booktabs,graphicx,xcolor}
\usepackage{hyperref}\hypersetup{colorlinks=true,linkcolor=blue!55!black,urlcolor=blue!55!black}
\usepackage{caption}\captionsetup{font=small,labelfont=bf}
\graphicspath{{figures/}}
\newcommand{\Prob}{\mathbb{P}}\newcommand{\Ex}{\mathbb{E}}
\newcommand{\A}[1]{\texttt{ALLOCATION\_#1}}
\title{\textbf{A mean-field theory for the blackboard game}\\[2mm]
\large Derivation, solution, and an honest comparison with the exact simulation}
\author{MA-CC blackboard-control study}\date{4 October 2026}
\begin{document}\maketitle

\begin{abstract}\noindent
We derive a mean-field theory for the blackboard game from its microscopic
rules, with \textbf{no fitted parameter}: every coefficient is a game setting
($\rho$, $N$, $q$, $b$). We solve it numerically and compare it with the exact
simulation. \textbf{It does not match}: mean absolute error $0.25$ in the
population coordinate, and in the controlled arms the theory gets the
\emph{direction} wrong. We identify the cause precisely --- a product-form
closure cannot represent a posterior whose value depends on whether particular
facts \emph{co-occur} --- and state what the theory needs. We also give the
continuous-time limit and the Ornstein--Uhlenbeck reduction that the closure
would permit if it worked.
\end{abstract}

\section{What to track, and why}

The game's only epistemic dynamics are \emph{decay} and
\emph{exposure-activation}, and both act on an (agent, fact) pair. So the
natural mean-field variable is the occupancy
\[
m_{if}(t) \;=\; \Prob\big(\text{fact } f \text{ is active for agent } i \text{ at round } t\big).
\]
The population coordinate we measure follows by averaging the exact posterior,
\begin{equation}
M(t) \;=\; \frac1N\sum_i \Ex_{K_i}\big[\Prob(\A{0}\mid K_i)\big],
\qquad K_i \sim \textstyle\prod_f \mathrm{Bernoulli}(m_{if}),
\label{eq:closure}
\end{equation}
and Eq.~\eqref{eq:closure} is the \textbf{closure}: it asserts that an agent's
facts are independent given their marginals. Everything that follows is exact
except this.

\section{Derivation, one round}

In the order the runtime applies them.

\paragraph{1. Decay.} Each active fact survives with probability $\rho$:
\begin{equation}
m_{if} \;\longmapsto\; \rho\, m_{if}.
\end{equation}

\paragraph{2. The board.} Agent $i$ posts one fact it holds. Writing $\pi_{if}$
for the probability it chooses $f$, and approximating the choice as uniform over
held facts, $\pi_{if} = m_{if}/\sum_g m_{ig}$. The controller adds its $b$ facts,
indicated by $c_f \in \{0,1\}$. The per-post marginal is
$\bar\pi_f = \frac1N\sum_i \pi_{if}$.

\paragraph{3. Exposure.} This is where the microscopic detail matters. Agents
are updated \emph{sequentially} and the board \emph{builds up within the round},
so the agent in position $k$ reads only the $k$ posts made before it. With
per-post probability $\bar\pi_f$, the chance that position $k$ sees $f$ is
$1-(1-\bar\pi_f)^k$, and averaging over positions gives a closed form:
\begin{equation}
\bar u_f \;=\; \frac1N\sum_{k=0}^{N-1}\Big[1-(1-\bar\pi_f)^k\Big]
\;=\; 1-\frac{1-(1-\bar\pi_f)^N}{N\,\bar\pi_f}.
\label{eq:seq}
\end{equation}
A \emph{synchronous} mean field would instead use $u_f = \mathbf 1\{B_f\ge 1\}$,
which with full board reading sends every occupancy to one after a single round.
Equation~\eqref{eq:seq} is the damping that replaces it. A controller post is on
the board before any agent reads, so it is seen by everyone:
$u_f = 1-(1-\bar u_f)(1-c_f)$.

Under \texttt{uniform} sampling of $q$ of $n$ messages, the hypergeometric form
replaces Eq.~\eqref{eq:seq},
$u_f = 1-\binom{n-k_f}{q}\big/\binom{n}{q}$, which is exactly Eq.~(7) of the
task\_004 paper.

\paragraph{4. Combine.} A fact is active afterwards if it survived decay or was
(re)activated:
\begin{equation}
\boxed{\;m_{if}(t+1) \;=\; \rho\,m_{if}(t) \;+\; \big(1-\rho\,m_{if}(t)\big)\,u_f(t)\;}
\label{eq:map}
\end{equation}

\paragraph{What it depends on.} Only $\rho$, $N$, $q$, $b$, the controller rule,
and the initial assignment. \textbf{Nothing is fitted.} That is the theory's
strongest claim and it is what makes the comparison below a real test.

\section{Continuous time, and the Ornstein--Uhlenbeck reduction}

Writing $\rho=e^{-\gamma}$ and taking one round as $\mathrm dt$,
Eq.~\eqref{eq:map} becomes
\begin{equation}
\dot m_{if} \;=\; -\gamma\, m_{if} \;+\; \big(1-m_{if}\big)\,\alpha_f,
\qquad \gamma=-\ln\rho,\quad \alpha_f=\text{exposure rate},
\label{eq:cont}
\end{equation}
a logistic relaxation with fixed point $m^\ast_f=\alpha_f/(\gamma+\alpha_f)$ and
rate $\gamma+\alpha_f$. This is the discrete-to-continuous correspondence: the
\textbf{deactivation rate is $-\ln\rho$}, exactly the $\gamma$ of the task\_004
paper's Table~1, and the per-round map is its Euler step. The two descriptions
agree to $O(\mathrm dt^2)$ and diverge when $\gamma\,\mathrm dt$ is not small,
i.e.\ at low $\rho$: at $\rho=0.6$, $\gamma=0.51$ per round, so a one-round step
is already a $50\%$ correction and the discrete map should be preferred.

If the closure held, linearising Eq.~\eqref{eq:cont} about $m^\ast$ and adding
finite-$N$ sampling noise would give the three-variable Ornstein--Uhlenbeck
system requested earlier,
\[
\mathrm dM=-\theta(M-M_\infty)\mathrm dt+g_p S^p\,\mathrm dt+g_c S^c\,\mathrm dt+\sigma_M\mathrm dW,
\quad
\mathrm dS^p=-\theta_p(S^p-M)\mathrm dt+\sigma_p\mathrm dW_p,
\]
\[
\mathrm dS^c=-\theta_c\big(S^c-u\,\pi(M)\big)\mathrm dt+\sigma_c\mathrm dW_c,
\]
with $\theta=-\ln\rho$, $M_\infty=1/3$, the peer force reverting to $M$ itself,
and the controller force reverting to the activation policy evaluated at $M$ ---
the feedback. Its stationary covariance would follow from the Lyapunov equation
and the path KL in closed form from Girsanov. \textbf{We do not report those
numbers}, because the closure fails first.

\section{The comparison, and it fails}

\begin{figure}[h]\centering\includegraphics[width=\textwidth]{fig_theory}
\caption{Theory (dashed) against the exact simulation (solid), no fitting.
Silent arms in grey, false-control arms in red.}
\end{figure}

\begin{center}\small
\begin{tabular}{llrrr}\toprule
setup & arm & sim $e$ (r29) & theory $M$ & error\\\midrule
nosolution-v2 $\rho{=}1$ & silent & 0.779 & 0.500 & $-0.279$\\
nosolution-v2 $\rho{=}1$ & false & 0.125 & 0.500 & $+0.375$\\
nosolution-v2 $\rho{=}0.75$ & silent & 0.545 & 0.504 & $-0.041$\\
symmetric-v2 $\rho{=}0.75$ & silent & 0.598 & 0.759 & $+0.161$\\
symmetric-v2 $\rho{=}0.75$ & false & 0.067 & 0.654 & $+0.587$\\
\bottomrule\end{tabular}\end{center}

\noindent
Mean absolute error $0.252$, maximum $0.587$. The silent arms are sometimes
close (symmetric-v2 at $\rho=0.75$ agrees to $0.03$ over rounds 2--10); the
\textbf{controlled arms are wrong in direction}, the theory rising where the
simulation falls.

\section{Why it fails, precisely}

\paragraph{The closure, not the dynamics.} The occupancy map
\eqref{eq:map} is close to exact. The failure is Eq.~\eqref{eq:closure}: the
posterior $\Prob(\A{0}\mid K)$ is a strongly nonlinear function of \emph{which}
facts co-occur. In symmetric-v2 the six decisive facts prove \A{0} only
\emph{together}; a product measure with the right marginals almost never
realises that conjunction, and almost never realises the conjunctions that
make \A{2} favourite either. Averaging over independent facts therefore pulls
every agent toward a generic mid-range posterior --- which is exactly the
$M\approx0.5$--$0.76$ plateau the theory shows.

\paragraph{Consequence for the controlled arms.} The simulated oracle drives $e$
to $0.067$ by placing \emph{a few specific} \A{2}-favouring facts in everyone's
hands. Under a product measure those facts are diluted among the agent's other
facts and their joint effect never forms. The theory consequently sees the
controller as adding evidence in general, which raises $M$.

\paragraph{What to change.} In increasing order of cost:
\begin{enumerate}
\item \textbf{Keep the exact conjunction structure for a small subset.} Track the
joint law of the 6 decisive facts (64 states) exactly and the rest by marginals.
Cheap, and it is where the nonlinearity lives.
\item \textbf{Pair correlations.} Add $m_{ifg}$ for pairs, closing at third
order. $\sim\!20^2/2$ numbers per agent; still trivial to integrate.
\item \textbf{Change the coarse variable.} Track not $\Ex[\Prob(\A{0}|K)]$ but
the distribution of \emph{proof coverage} --- how many of each agent's minimal
proofs are complete. That variable is closer to additive in the facts.
\end{enumerate}

\paragraph{Does the theory need fitting?} \textbf{No, and it should not be
fitted.} Every coefficient is a game setting. If a fitted parameter were
introduced to close the $0.25$ gap it would absorb the closure error and the
theory would stop being a test of anything. The honest statement is that the
product-form mean field is \emph{falsified} for this observable, by a
simulation that costs 28 seconds to run.

\section{What this exercise established}

The mean-field programme is not dead; its first closure is. The occupancy
dynamics \eqref{eq:map}, the sequential-exposure correction \eqref{eq:seq} and
the identification $\gamma=-\ln\rho$ are all reusable. What fails is the
assumption that an agent's facts are independent --- and that assumption is
exactly what the game is designed to violate, since the whole point of a
\emph{decisive set} is that certain facts matter only in conjunction.

That is a result worth having before any of it reaches a paper, and it was
obtained for free.

\end{document}
